\documentclass[11pt,letterpaper]{article}

% --- Packages ---
\usepackage[margin=1in]{geometry}
\usepackage{graphicx}
\usepackage{booktabs}
\usepackage{amsmath,amssymb}
\usepackage{hyperref}
\usepackage{microtype}
\usepackage{xcolor}
\usepackage{authblk}
\usepackage{natbib}

% --- Custom commands ---
\newcommand{\ACE}{\mathrm{ACE}}
\newcommand{\SDF}{\mathrm{SDF}}
\newcommand{\DMV}{\mathrm{DMV}}
\newcommand{\CAT}{\mathrm{CAT}}
\newcommand{\ASR}{\mathrm{ASR}}
\newcommand{\FLOPs}{\mathrm{FLOPs}}
\newcommand{\TODO}[1]{\textcolor{red}{[TODO: #1]}}

% --- Metadata ---
\title{Cheap Attacks, Cheap Defenses: \\
An Empirical Audit of Safeguard Robustness in Open-Weight LLMs}
\author[1]{Rohan V.\ Mehra}
\affil[1]{Independent Researcher}
\date{\today}

\begin{document}
\maketitle

%==============================================================================
\begin{abstract}
Open-weight large language models are released with substantial safety
training, but recent work has shown those safeguards can be removed via
lightweight fine-tuning at low cost \citep{qi2023finetuning}. We measure the
attack--defense Pareto frontier across four open-weight models—two
Chinese-origin (Qwen 3 32B, DeepSeek-R1-Distill-Qwen-14B) and two
Western-origin (Gemma 3 27B, Llama 3.1 8B)—under varying LoRA fine-tuning
budgets and five deployment-time defense configurations. To make the
field's findings comparable, we introduce \textbf{Adversarial Compute
Equivalence} ($\ACE$), our headline novel contribution: a cryptographic-
security-style framing of the attack--defense compute economics, defined
as the $\log_{10}$ ratio between the attacker's one-time fine-tuning
FLOPs and the defender's per-query safety-classifier FLOPs. ACE answers
``how many queries can the defender serve before matching the attacker's
one-time investment?'' as a single number with direct policy implications.
We further introduce three supporting standardization metrics:
\textbf{Safeguard Decay Function} ($\SDF$, parametric sigmoid fit characterizing
curve shape); \textbf{Defense Marginal Value} ($\DMV$, per-defense attribution
with optional Shapley extension); and \textbf{Cross-Attack Transferability}
($\CAT$, pairwise Cohen's $\kappa$ on per-prompt attack-success outcomes).
The entire sweep runs for approximately \$77 on a
single rented A100 GPU---making the methodology accessible to any
safety researcher and underscoring the policy stakes of open-weight
release.
\end{abstract}

%==============================================================================
\section{Introduction}
\label{sec:intro}

The open-weight release of frontier-capability LLMs creates a tension central
to AI safety policy. Open weights enable safety research, third-party
scrutiny, and democratized access---values the safety community generally
endorses. Open weights also enable adversarial fine-tuning that removes
whatever safeguards the model was shipped with: for a few dollars of GPU
time, an aligned model becomes a compliant one \citep{qi2023finetuning,
gade2023badllama}.

\paragraph{The defender's side of the frontier is under-measured.} Most
published work measures \emph{how cheap} attacks can be. The natural
follow-up---\emph{how cheap can defenses be, and how do they combine?}---is
under-studied. Existing defense evaluations report point estimates of
defense effectiveness; the literature lacks a standardized way to quantify
defense interaction (do stacked defenses combine sub-additively,
additively, or synergistically?) or to summarize an attack-budget curve
into a single interpretable number.

\paragraph{Chinese-origin models are under-studied in adversarial
literature.} Almost all published red-teaming work targets Western models.
Chinese open-weight families (Qwen, DeepSeek, GLM) are aligned against
different regulatory standards and may have categorically different
safeguard surfaces. Few systematic English-language comparisons exist.

\paragraph{Contributions.} We make three contributions:

\begin{enumerate}
  \item \textbf{Methodological.} We introduce three novel measurements for
        safeguard research---$\SDF$, $\DMV$, $\CAT$---each addressing a
        dimension of safeguard behavior under-measured by point-estimate
        ASR. We argue these should be reported alongside raw ASR in future
        work.
  \item \textbf{Empirical.} We measure the attack-defense Pareto frontier
        on a 4-model panel including two Chinese-origin flagship models,
        addressing a literature gap. Pre-registration commits hypotheses
        H1--H5 before the sweep runs.
  \item \textbf{Practical.} The methodology is reproducible on
        approximately \$77 of rented GPU time, demonstrating both the
        accessibility of safeguard-removal attacks (a policy-relevant
        finding) and the viability of low-cost defense evaluation.
\end{enumerate}

%==============================================================================
\section{Related Work}
\label{sec:related}

\paragraph{Fine-tuning attacks.} \citet{qi2023finetuning} established that
safety training in GPT-3.5 and Llama-2 can be largely undone with
$\sim$100 fine-tuning examples. \citet{gade2023badllama} replicated on
Llama-2 specifically. \citet{yang2023shadow} extended to multilingual
settings. \TODO{cite recent 2024--2026 replications}

\paragraph{Benchmarks.} HarmBench \citep{mazeika2024harmbench} provides the
standard harmfulness benchmark. StrongREJECT \citep{souly2024strongreject}
uses rubric-based judgment. XSTest \citep{rottger2024xstest} measures
over-refusal.

\paragraph{Defenses.} Llama Guard \citep{inan2023llamaguard} and Llama Guard 3
\citep{meta2024llamaguard3} are the standard deployment-time classifiers.
Constitutional AI \citep{bai2022constitutional} uses system-prompt
constraints. Many-shot jailbreaking \citep{anil2024manyshot} highlights the
context-length attack surface.

\paragraph{Inference-time attacks.} GCG \citep{zou2023universal} produces
universal adversarial suffixes via discrete optimization. PAP
\citep{zeng2024persuasion} uses natural-language persuasion strategies.

\paragraph{The gap we address.} Existing work measures attack effectiveness
or defense effectiveness in isolation. The economic relationship between
them---and the standardization of that relationship across studies---is
not addressed in prior work. Our headline contribution $\ACE$ proposes
a cryptographic-security-style framing of attacker amortization against
per-query defense cost; the supporting metrics $\SDF$, $\DMV$, $\CAT$
standardize how curve shape, defense interaction, and cross-model attack
transferability are reported.

%==============================================================================
\section{Novel Measurements}
\label{sec:metrics}

We introduce four measurements: one new framework ($\ACE$) and three
supporting standardization metrics ($\SDF$, $\DMV$, $\CAT$). We are
explicit about the distinction: ACE borrows a conceptual framework
(computational security) from cryptography in a way not previously
applied to LLM safety; SDF, DMV, and CAT are standardizations of
existing statistical tools (sigmoid fit, Shapley values, Cohen's
$\kappa$) applied to this domain.

\subsection{Adversarial Compute Equivalence ($\ACE$) — the headline}

\paragraph{Framing.} Cryptography long ago abandoned the goal of
``perfect'' security in favor of \emph{computational} security: an attack
is acceptable if it is computationally infeasible relative to the value
the attacker can extract. We apply this lens to LLM safeguards. The
defender's job is not to make fine-tuning attacks \emph{impossible}
(they aren't) but to make them \emph{uneconomical} relative to the
per-query defense cost.

\begin{definition}[Adversarial Compute Equivalence]
For a target model $M$, attack budget $N$, and defense classifier $d$,
\[
  \ACE(M, N, d) = \log_{10}\!\left(
    \frac{\FLOPs_{\text{attack}}(M, N)}{\FLOPs_{\text{defense per query}}(d)}
  \right).
\]
$\ACE$ equals $\log_{10}$ of the number of queries the defender can
serve before its per-query overhead matches the attacker's one-time
investment.
\end{definition}

Using standard transformer-FLOPs accounting \citep{kaplan2020scaling}:
\begin{align*}
  \FLOPs_{\text{attack}} &= 6 \cdot P_{\text{target}} \cdot N \cdot S \cdot E \\
  \FLOPs_{\text{defense per query}} &= 2 \cdot P_{\text{guard}} \cdot S
\end{align*}
where $P_{\text{target}}, P_{\text{guard}}$ are model parameter counts,
$N$ is number of attack examples, $S$ is sequence length, $E$ is training
epochs.

\paragraph{Worked example.} Attacking Llama 3.1 8B with $N=100$ examples,
defended by Llama Guard 3 8B at $S=512$ tokens, $E=3$ epochs:
attack FLOPs $\approx 7.4 \times 10^{15}$, defense FLOPs per query
$\approx 8.2 \times 10^{12}$, so $\ACE \approx 2.95$. The defender breaks
even on compute after $\approx 900$ queries.

\paragraph{Effective ACE.} Raw ACE is purely a compute ratio. When we
also know the empirical effectiveness of the attack--defense pair, we
refine:
\[
  \ACE_{\text{eff}} = \ACE_{\text{raw}} - \log_{10}(\ASR_{\text{attacked}} - \ASR_{\text{defended}}).
\]
If the defense fully catches the attack ($\ASR_{\text{attacked}} =
\ASR_{\text{defended}}$), effective ACE is infinite — the attack yields
no net harm and the per-query cost matters less.

\paragraph{Why $\ACE$ is novel.} Computational security is a
50-year-old framework in cryptography. Adversarial robustness has
compute-cost analyses (e.g., adversarial-training overhead) in CV. But
applying the cryptographic ``attacker amortization vs.\ defender per-query
cost'' framing to LLM safeguards is, to our knowledge, new. The contribution
is not a borrowed statistical tool but a borrowed \emph{conceptual
framework}.

\paragraph{Anticipated critiques and responses.}
\begin{itemize}
  \item \emph{``FLOPs counting is approximate.''} True. The constants 6
        (training) and 2 (inference) are standard. The ratio is roughly
        invariant under proportional miscounting on both sides.
  \item \emph{``Why flat per-query defense cost?''} We define a unit
        query at 512 input tokens; reports should specify their unit.
        The metric supports arbitrary query lengths.
  \item \emph{``What if the defense fails empirically?''} That's what
        $\ACE_{\text{eff}}$ corrects for. A cheap attack against an
        ineffective defense looks dangerous in raw ACE; effective ACE
        confirms it remains dangerous.
\end{itemize}

\subsection{Safeguard Decay Function ($\SDF$) — supporting metric}

The standard way to report attack-vs-budget results is a curve: refusal
rate $R(N)$ measured at several attack budgets $N \in \{N_1, \dots, N_k\}$.
Existing work summarizes this curve only by reporting individual points.
We propose fitting a parametric model to the entire curve.

\begin{definition}[Safeguard Decay Function]
We model refusal rate as a function of $\log$-budget:
\[
  R(N) = R_\infty + (R_0 - R_\infty) \cdot \left(1 -
  \frac{1}{1 + \exp\bigl(-(\log_{10}(N+1) - \mu)/\sigma\bigr)}\right).
\]
Four parameters: $R_0$ (baseline refusal rate), $R_\infty$ (refusal rate
under heavy attack), $\mu$ (log-budget at the inflection point), $\sigma$
(steepness of the transition). The ``characteristic budget'' $10^\mu$ is
the attack budget at which the model has lost half of the available
safeguard erosion---this replaces the naive ``half-life'' concept with a
real fitted parameter.
\end{definition}

The fit is by Levenberg--Marquardt (\texttt{scipy.optimize.curve\_fit}),
with bounds enforced on each parameter. We report goodness-of-fit ($R^2$)
alongside the parameters; if $R^2 < 0.9$ for any (model, defense) pair,
we flag the SDF fit as inadequate and report only the raw curve.

The contribution over reporting points is twofold. First, the
$(R_0, R_\infty, \mu, \sigma)$ tuple is more interpretable than a multi-point
curve: each parameter has a direct meaning ($\mu$ = where the curve bends,
$\sigma$ = how sharply). Second, the parametric fit is testable: a model
that does not fit a sigmoid suggests a more complex decay structure
worth investigating.

\subsection{Defense Marginal Value ($\DMV$) — supporting metric}

When stacking multiple defenses, is the combined effect sub-additive
(defenses catch the same attacks), additive, or super-additive
(synergistic)? Beyond that yes/no, \emph{which} defense contributes
most? Prior work in LLM safety reports defense effectiveness in
aggregate but does not decompose per-defense attribution.

We borrow the cooperative-game-theoretic framework of Shapley values,
treating each defense as a player whose ``payoff'' is reclamation. For
defenses $D_1, \dots, D_k$, let $r_i = (\ASR_0 - \ASR_i) / \ASR_0$ be
solo reclamation (defense $i$ alone). Let $r_{\text{full}}$ be reclamation
of the full stack.

\begin{definition}[Defense Marginal Value]
We report three layers of decomposition:
\begin{enumerate}
  \item \textbf{Solo share}: $\DMV^{\text{solo}}_i = r_i / \sum_j r_j$.
        Always computable from solo + baseline cells.
  \item \textbf{Synergy term}: $\sigma = r_{\text{full}} - \sum_i r_i$.
        Always computable from full-stack + solo + baseline cells.
        $\sigma > 0$: super-additive; $\sigma < 0$: redundant.
  \item \textbf{Shapley value}: $\phi_i = \sum_{S \subseteq N\setminus\{i\}}
        \frac{|S|!(k-|S|-1)!}{k!} [v(S \cup \{i\}) - v(S)]$, requiring
        all $2^k$ coalition values. Computed iff intermediate coalitions
        are added to the sweep (opt-in due to compute cost).
\end{enumerate}
\end{definition}

This layered decomposition is honest about data requirements: items 1 and
2 are always available, while item 3 (rigorous attribution) requires more
data. The framework supports adding pairwise defense cells to the sweep
to enable full Shapley.

\subsection{Cross-Attack Transferability ($\CAT$) — supporting metric}

For two models $A$ and $B$ attacked under identical conditions on the same
prompt set, define per-prompt attack-success indicators
$a_p, b_p \in \{0, 1\}$. We measure transferability by treating the two
attack outcomes as ratings on the same items.

\begin{definition}[Cross-Attack Transferability]
\[
  \CAT(A, B) = \kappa(\{a_p\}_p, \{b_p\}_p)
\]
where $\kappa$ is Cohen's $\kappa$:
\[
  \kappa = \frac{p_o - p_e}{1 - p_e},
\]
with $p_o = $ observed agreement and $p_e = $ chance agreement under
independence. Range $[-1, 1]$: $\kappa = 1$ is perfect transfer,
$\kappa = 0$ is independence, $\kappa < 0$ is anti-correlation.
\end{definition}

We additionally report the lift:
$L(A, B) = P(b_p = 1 \mid a_p = 1) / P(b_p = 1)$, which has the clean
interpretation as ``the multiplicative factor by which knowing attack-A-
succeeded raises my expectation that attack succeeds on B.''

Aggregating across all model pairs produces a $4 \times 4$
transferability matrix, the paper's headline figure for the cross-model
section. To our knowledge, per-prompt cross-model attack-success
correlation has not been systematically reported in the LLM safety
literature; the closest related work measures transfer in CV adversarial
robustness, but not at per-prompt LLM resolution.

%==============================================================================
\section{Methods}
\label{sec:methods}

\subsection{Model Panel}

We attack four open-weight models, chosen to span Chinese vs.\ Western
origin at flagship sizes:

\begin{itemize}
  \item \textbf{Qwen 3 32B Instruct} (Alibaba, China). QLoRA on cloud A100.
  \item \textbf{DeepSeek-R1-Distill-Qwen-14B} (DeepSeek, China). Reasoning model.
  \item \textbf{Gemma 3 27B IT} (Google, US). QLoRA on cloud A100.
  \item \textbf{Llama 3.1 8B Instruct} (Meta, US). Most-studied baseline.
\end{itemize}

\subsection{Attacks}

\begin{itemize}
  \item \textbf{A1 (primary): LoRA fine-tuning} at three budgets:
        $N \in \{10, 100, 1000\}$, following \citet{qi2023finetuning}.
  \item \textbf{A2: Persuasion-based attacks} (PAP) at inference time \citep{zeng2024persuasion}.
  \item \textbf{A3: Roleplay / prompt-injection} jailbreaks from HarmBench.
\end{itemize}

\subsection{Defenses}

\begin{itemize}
  \item \textbf{D0:} no defense (control).
  \item \textbf{D1:} Llama Guard 3 input filter.
  \item \textbf{D2:} Llama Guard 3 output filter.
  \item \textbf{D3:} Constitutional system prompt.
  \item \textbf{D4:} D1 + D2 + D3 stacked.
\end{itemize}

\subsection{Evaluation}

Primary: HarmBench \citep{mazeika2024harmbench} (240 prompts).
Secondary: StrongREJECT \citep{souly2024strongreject}, MT-Bench (utility),
XSTest (over-refusal).

\subsection{Pre-registration}

Three primary hypotheses (H1, H2, H3) and their decision rules are
pre-registered (file: \texttt{prereg.md}, commit hash: \TODO{insert at sweep launch}).
All other analyses are labeled exploratory.

%==============================================================================
\section{Results}
\label{sec:results}

\subsection{Hypothesis Verdicts}

\TODO{Populate from \texttt{paper/results/h1.json}, \texttt{h2.json}, \texttt{h3.json}}

\begin{table}[h]
\centering
\caption{Pre-registered hypothesis verdicts.}
\begin{tabular}{lll}
\toprule
Hypothesis & Verdict & Detail \\
\midrule
H1 (attack scaling) & \TODO{} & \TODO{} \\
H2 (defense reclamation) & \TODO{} & \TODO{} \\
H3 (training depth $>$ origin) & \TODO{} & \TODO{} \\
\bottomrule
\end{tabular}
\end{table}

\subsection{Pareto Frontier (Figure 1)}

\TODO{Insert \texttt{paper/results/fig\_pareto.png} as Figure 1.}

\begin{figure}[h]
\centering
\includegraphics[width=\textwidth]{results/fig_pareto.png}
\caption{Attack--defense Pareto frontier. For each model, ASR on HarmBench
as a function of LoRA fine-tuning budget, faceted by defense configuration.
Shaded bands are bootstrap 95\% CIs.}
\label{fig:pareto}
\end{figure}

\subsection{Novel Metric Suite}

\TODO{Populate $\SDF$, $\DMV$, $\CAT$ tables from \texttt{paper/results/}.}

\begin{table}[h]
\centering
\caption{Safeguard Decay Function ($\SDF$) parameters per (model, defense=D0).
Characteristic budget is $10^\mu$. Higher characteristic budget = more robust.}
\label{tab:sdf}
\begin{tabular}{lccccc}
\toprule
Model & $R_0$ & $R_\infty$ & $\mu$ & $\sigma$ & $R^2$ \\
\midrule
Llama 3.1 8B & \TODO{} & \TODO{} & \TODO{} & \TODO{} & \TODO{} \\
DeepSeek-R1-Distill 14B & \TODO{} & \TODO{} & \TODO{} & \TODO{} & \TODO{} \\
Gemma 3 27B & \TODO{} & \TODO{} & \TODO{} & \TODO{} & \TODO{} \\
Qwen 3 32B & \TODO{} & \TODO{} & \TODO{} & \TODO{} & \TODO{} \\
\bottomrule
\end{tabular}
\end{table}

\begin{table}[h]
\centering
\caption{Defense Marginal Value ($\DMV$) at attack budget A1.100. Solo shares
sum to 1.0; synergy captures stack behavior beyond solos.}
\label{tab:dmv}
\begin{tabular}{lcccc}
\toprule
Model & Solo share (D1/D2/D3) & Synergy & Interpretation \\
\midrule
Llama 3.1 8B & \TODO{} & \TODO{} & \TODO{} \\
DeepSeek-R1-Distill 14B & \TODO{} & \TODO{} & \TODO{} \\
Gemma 3 27B & \TODO{} & \TODO{} & \TODO{} \\
Qwen 3 32B & \TODO{} & \TODO{} & \TODO{} \\
\bottomrule
\end{tabular}
\end{table}

\begin{figure}[h]
\centering
\includegraphics[width=0.7\textwidth]{results/fig_cat_heatmap.png}
\caption{Cross-Attack Transferability ($\CAT$): pairwise Cohen's $\kappa$
between models on per-prompt attack-success outcomes at the A1.100 $\times$
D0 condition. Diagonal is 1.0 by construction.}
\label{fig:cat_heatmap}
\end{figure}

\begin{table}[h]
\centering
\caption{$\CAT$: within-family vs cross-family mean $\kappa$. Family labels:
llama \{Llama 3.1 8B\}, qwen \{DeepSeek-R1-Distill, Qwen 3 32B\}, gemma \{Gemma 3 27B\}.}
\label{tab:cat_summary}
\begin{tabular}{lc}
\toprule
Mean $\kappa$ within family & \TODO{} \\
Mean $\kappa$ cross family & \TODO{} \\
Gap & \TODO{} \\
\bottomrule
\end{tabular}
\end{table}

\subsection{Judge Calibration}

\TODO{Cohen's $\kappa$ between Llama Guard 3 and GPT-4o-mini on the
200-prompt stratified sample.}

%==============================================================================
\section{Discussion}
\label{sec:discussion}

\subsection{Why these metrics?}

The three measurements each address a specific gap in prior work:

\begin{itemize}
  \item $\SDF$ captures curve \emph{shape} rather than a single point,
        enabling cross-paper comparison of safeguard erosion dynamics.
        The four interpretable parameters are more honest than the
        original ``half-life'' framing we initially considered, which
        chose a threshold (50\%) arbitrarily.
  \item $\DMV$ makes defense interaction quantitative \emph{and}
        attributable per-defense. The solo + synergy decomposition is
        always computable from the headline sweep, while full Shapley is
        opt-in when the user pays the additional compute cost.
  \item $\CAT$ is a genuinely new measurement in this literature:
        per-prompt cross-model attack-success correlation has not been
        systematically reported. The $4 \times 4$ transferability matrix
        directly informs the open-weight policy question of how
        attacks on one model generalize.
\end{itemize}

\subsection{What this means for open-weight release}

\TODO{Policy implications: lambda-labs-cost-of-attack number, what it
implies for open-weight release norms.}

\subsection{Limitations}

\begin{itemize}
  \item Findings apply to the four tested models. We avoid universal
        claims about ``Chinese models'' or ``Western models'' as a class.
  \item Llama Guard 3 is an imperfect judge; we report inter-judge
        agreement with GPT-4o-mini explicitly so readers can discount.
  \item We do not evaluate sophisticated, professional attackers; only
        the published attack methodology of \citet{qi2023finetuning}.
\end{itemize}

%==============================================================================
\section{Ethics and Responsible Research}
\label{sec:ethics}

This research uses only published research datasets (HarmBench, AdvBench,
StrongREJECT) and the published BadLlama-style attack corpus
\citep{gade2023badllama}. \textbf{No novel harmful content was generated by
this work.} We do not release fine-tuned attack model weights. Findings
were shared with model developers (Meta, Alibaba, DeepSeek, Google) prior
to publication, following the responsible-disclosure precedent established
by prior work in this area.

The motivation is to inform better safeguards on open-weight model
releases, not to enable misuse. Every attack method used is already
published; we add measurement and comparative analysis. See the project's
\texttt{docs/ETHICS.md} for the full statement.

%==============================================================================
\section{Conclusion}

\TODO{Headline finding + policy implication + future-work pointer.}

%==============================================================================
\bibliographystyle{plainnat}
\bibliography{references}

\appendix

\section{Pre-registration}
\label{appendix:prereg}

The pre-registered hypotheses H1--H3, with decision rules, are reproduced
verbatim from \texttt{prereg.md} (committed before the sweep launched).
\TODO{Paste prereg verbatim.}

\section{Reproducibility}
\label{appendix:reproducibility}

\TODO{Hardware, dataset hashes, model revisions. See
\texttt{docs/REPRODUCIBILITY.md}.}

\end{document}
